\exnumber{05}
\extitle{UX II - Ex 6 - Parcours utilisateur}
\exfiles{ex05/archive zip, nomme Name\_FirstName\_DeliverableName\_Date\_VersionNumber}
\exauthorize{Ceux nommés dans l'exercice}
\makeheaderfiles

\textbf{Instructions :}

Structurez vos idées principales en un user flow complet.
Vous pouvez inclure une, deux ou les trois idées principales de votre processus d'idéation, selon ce qui serait le plus cohérent pour les utilisateurs.

\textbf{Livrable :}

\begin{itemize}
\item Un parcours utilisateur complet sur un board (numérique ou photos d'une version papier)
\item Déposez les documents sur Github, dans un zip nommé « ProjectTitle\_Name\_FirstName », et chaque livrable ainsi : « Name\_FirstName\_DeliverableName\_Date\_VersionNumber ». Exemple : « Doe\_John\_MidFiPrototype\_05242024\_V1 ».
\item Ajoutez un titre dans tous les documents écrits.
\item Soignez l'aspect global des livrables. Ils doivent être clairs, simples et lisibles. Vous pouvez regarder canva.com, sans surcharger le dessin de détails.
\end{itemize}

\textbf{Ressources :}

\begin{itemize}
\item Glossaire User flow \url{https://www.productplan.com/glossary/user-flow/}
\item Site flows vs User Flows \url{https://uxmovement.com/wireframes/site-flows-vs-user-flows-when-to-use-which/}
\item Comment réaliser un diagramme de User Flow \url{https://www.lucidchart.com/blog/how-to-make-a-user-flow-diagram}
\end{itemize}

\textbf{Citation :}

« Chaque cas d'usage est représenté comme une séquence d'étapes simples, commençant par l'objectif d'un utilisateur et se terminant lorsque cet objectif est atteint. » Usability.gov

